% ===== PRÉAMBULE CANONIQUE — à coller VERBATIM en tête de CHAQUE .tex =====
\documentclass[11pt,a4paper]{article}
\usepackage[utf8]{inputenc}\usepackage[T1]{fontenc}\usepackage{lmodern}
\usepackage[french]{babel}
\usepackage{amsmath,amssymb,mathtools}\usepackage{icomma}
\usepackage{siunitx}\sisetup{locale=FR}  % \qty{}{}, \si{}, \ang{}
\usepackage{xcolor}\usepackage{colortbl}
\usepackage{tikz}\usepackage{pgfplots}\pgfplotsset{compat=1.18}
\usepackage[siunitx,europeanresistors]{circuitikz} % circuits (élec.)
\usepackage[most]{tcolorbox}\usepackage{enumitem}\usepackage{array,booktabs}
\usepackage[a4paper,margin=2.2cm]{geometry}
\usetikzlibrary{arrows.meta,calc,angles,quotes,positioning,patterns,%
  backgrounds,shapes.geometric,decorations.pathmorphing,decorations.markings,intersections}
\definecolor{primary}{RGB}{222,49,99}\definecolor{primarydark}{RGB}{150,22,55}
\definecolor{defcol}{RGB}{0,114,178}\definecolor{methcol}{RGB}{52,92,148}
\definecolor{excol}{RGB}{0,135,90}\definecolor{attncol}{RGB}{200,40,55}
\tcbset{boxbase/.style={enhanced,breakable,boxrule=0.9pt,arc=2.5pt,
  left=3mm,right=3mm,top=2.3mm,bottom=2.3mm,fonttitle=\bfseries,coltitle=white}}
% --- boîtes TUTORIEL (noms EXACTS, ne pas renommer) ---
\newtcolorbox{cle}[1][]{boxbase,colback=defcol!6,colframe=defcol,
  title={\textsc{À retenir}\ifx\relax#1\relax\relax\else\textmd{ : \itshape #1}\fi}}
\newtcolorbox{methode}[1][]{boxbase,colback=methcol!7,colframe=methcol,sharp corners,
  title={\textsc{Méthode}\ifx\relax#1\relax\relax\else\textmd{ --- \itshape #1}\fi}}
\newtcolorbox{exemplebox}[1][]{boxbase,colback=excol!5,colframe=excol,
  title={\textsc{Exemple}\ifx\relax#1\relax\relax\else\textmd{ : \itshape #1}\fi}}
\newtcolorbox{piege}[1][]{boxbase,colback=attncol!6,colframe=attncol,
  title={\textsc{Piège}\ifx\relax#1\relax\relax\else\textmd{ : \itshape #1}\fi}}
\newtcolorbox{remarque}[1][]{boxbase,colback=black!4,colframe=black!45,
  title={\textsc{Remarque}\ifx\relax#1\relax\relax\else\textmd{ : \itshape #1}\fi}}
% --- boîtes EXERCICE / CORRIGÉ (noms EXACTS, auto-comptées) ---
\newtcolorbox[auto counter]{exo}[1][]{boxbase,colback=primary!4,colframe=primary,
  title={\textsc{Exercice}~\thetcbcounter\ifx\relax#1\relax\relax\else\textmd{ --- \itshape #1}\fi}}
\newtcolorbox[auto counter]{corrige}[1][]{boxbase,colback=excol!4,colframe=excol,
  title={\textsc{Corrigé}~\thetcbcounter\ifx\relax#1\relax\relax\else\textmd{ --- \itshape #1}\fi}}
\newcommand{\vv}[1]{\vec{#1}}\newcommand{\ds}{\displaystyle}
\newcommand{\R}{\mathbb{R}}\newcommand{\N}{\mathbb{N}}\newcommand{\Z}{\mathbb{Z}}
% =========================================================================
\begin{document}
\sisetup{per-mode=symbol}
% ================= CORRIGÉ DIFFICILE =================

\begin{corrige}[la traction du câble]
\begin{enumerate}[label=\alph*)]
  \item Deux forces : la traction $T$ (vers le haut) et le poids $Mg$ (vers le bas).
        Axe vers le haut : $T-Mg=Ma$.
  \item $T=M(g+a)=800\times(10+1{,}5)=800\times 11{,}5=\qty{9200}{\newton}$.
  \item En MRU, $a=0$ : $T=Mg=800\times 10=\qty{8000}{\newton}$ (le câble équilibre
        juste le poids).
\end{enumerate}
\end{corrige}

\begin{corrige}[la corde en chute libre]
\begin{enumerate}[label=\alph*)]
  \item En chute libre, l'ensemble n'est soumis qu'à la pesanteur : $a=g$ (vers le bas).
  \item Sphère $B$, axe orienté vers le bas (sens du mouvement) :
        $m_B\,g-T=m_B\,a=m_B\,g$. On en tire $T=m_B g-m_B g=\qty{0}{\newton}$.
  \item \textbf{Non, la corde ne tire plus} ($T=0$). Au repos elle supportait
        $T=m_B g=1\times 10=\qty{10}{\newton}$ ; en chute libre, $A$ et $B$ tombent
        \emph{ensemble} avec la même accélération $g$, la corde devient inutile.
\end{enumerate}
\end{corrige}

\begin{corrige}[poids apparent --- vrai ou faux]
\begin{enumerate}[label=\alph*)]
  \item \textbf{Vrai.} MRU $\Rightarrow a=0\Rightarrow N=mg$.
  \item \textbf{Vrai.} $a>0\Rightarrow N=m(g+a)>mg$.
  \item \textbf{Vrai.} Chute libre $\Rightarrow a=-g\Rightarrow N=0$.
  \item \textbf{Faux.} Le poids \emph{réel} $mg$ ne change jamais (si $g$ est constant) :
        c'est le poids \emph{apparent} $N$ qui diminue.
\end{enumerate}
Trois propositions correctes (a, b, c).
\end{corrige}

\begin{corrige}[un trajet complet vers le haut]
On applique $N=m(g+a)$ à chaque phase.
\begin{enumerate}[label=\alph*)]
  \item $N=60\times(10+2)=\qty{720}{\newton}$ : elle se sent \textbf{plus lourde}.
  \item MRU ($a=0$) : $N=mg=\qty{600}{\newton}$ : sensation \textbf{normale}.
  \item $N=60\times(10-2)=\qty{480}{\newton}$ : elle se sent \textbf{plus légère}.
\end{enumerate}
On se sent plus lourd en accélérant vers le haut, plus léger en freinant à l'arrivée.
\end{corrige}

\end{document}
